\section{Historical context and scope of the comparison}
\label{app:context}

\paragraph{The useful lesson of Move 37.}
AlphaGo combined learned evaluation with tree search~\cite{alphago}.
After the 2016 match with Lee Sedol, Hassabis reported that its thirty-seventh
move in Game 2 had an estimated one-in-ten-thousand chance of being played by
a human~\cite{move37}. That estimate came from the system's account of human
play; it is not a universal probability of creativity. An improbable move can
be poor, and a valuable move can be unsurprising. InnovationZero consequently
uses Move42 as a prospective, explicitly defined research event, rather than
borrowing the historical estimate as a novelty threshold already satisfied by
its programs. The analogy motivates an experiment; it does not provide its
result.

\paragraph{Learning from search is an established idea.}
AlphaGo Zero removed the human-game supervision used by its predecessor,
training a policy and value network through search-guided self-play. Its
network improved search, which in turn supplied stronger training targets.
The paper's policy-only evaluation also distinguishes learned competence from
the strength obtained by adding search~\cite{alphagozero}. This is the relevant
precedent for the proposed feedback-budget curve. It does not make a
one-step contextual-bandit problem equivalent to a competitive, sequential
board game. Nor does learning without human game records imply learning
without human assumptions: a representation, a legal action space, an
objective, and a learning procedure still have designers.

\paragraph{From games to algorithms.}
AlphaTensor recasts matrix-multiplication algorithm discovery as a single-player
game; AlphaDev searches assembly programs using reinforcement learning and
measured performance~\cite{alphatensor,alphadev}. Their achievements show that
formal structure and executable feedback can support useful discoveries.
FunSearch and AlphaEvolve provide a complementary lineage: a language model
proposes code, evaluators assess it, and successful candidates inform later
proposals~\cite{funsearch,alphaevolve}. These systems are precedents for
outcome-driven discovery, rather than evidence that all invention necessarily
requires the same amount of search. Language models are participants in this
lineage; treating them as intrinsically incapable of contributing new
algorithms would misdescribe the evidence.

\paragraph{The closest technical antecedents.}
DeepCoder learns properties of programs from input-output examples to guide
symbolic search~\cite{deepcoder}. DreamCoder combines learned abstractions with
neural guidance over a growing program library~\cite{dreamcoder}.
AutoML-Zero evolves complete learning procedures from basic mathematical
operations~\cite{automlzero}. HyperNetworks generate the parameters of another
network~\cite{hypernetworks}; neural symbolic regression can predict a full
expression, including constants, before optional numerical
refinement~\cite{symbolicregression}. Latent Programmer predicts a discrete
program representation before generating source, organizing synthesis as a
two-level search~\cite{latentprogrammer}. HyperFast uses a meta-trained
hypernetwork to produce a dataset-specific classifier in a forward
pass~\cite{hyperfast}. These are close precedents for a solution representation
and data-conditioned generation, respectively; a fitted classifier and a
reusable learning procedure are different output objects. Together these works rule out a broad
priority claim for learned data-to-program generation, representation
learning over algorithms, or transferring computation from deployment to
training. The distinctive question here is narrower: can an outcome-trained
proposer commit a useful, reusable fitting procedure for a previously unseen
dataset family before receiving any candidate-specific feedback, and can that
boundary be audited independently of its later fit and evaluation?

\paragraph{The Eureka Machine frame.}
Socher's \emph{The Eureka Machine} connects scientific discovery to increasingly
automated cycles of hypothesis and experiment~\cite{socher}. In an
associated author-written account, he describes four foundations: accumulated
human knowledge, scientific measurements, digital simulations, and autonomous
physical laboratories~\cite{socherpillars}. Coordinating agents sit above these
foundations. This is a vision for a scientific system, not a claim that a
single model already implements every component. A data-conditioned algorithm
proposer could occupy one component of such a system: it attempts to transfer
lessons from earlier experiments into the first proposal for a new problem.
Its usefulness would not eliminate the need to test the resulting scientific
claim.

Recursive's research report describes multiple concurrent research threads
that propose, implement, run, and validate experiments before choosing what
to attempt next~\cite{recursiveresearch}. The company subsequently attributed
much of Flint's development stack to its Eureka Machine
system~\cite{recursiveflint}. These are company-reported development results;
they are neither independent validation of unrestricted recursive
self-improvement nor demonstrations of feedback-free scientific discovery.
The comparison clarifies where feedback is used. An iterative research system
can produce training evidence for an amortized proposer, while a strong
proposer can improve the starting point of an iterative research system.
Recent work on SwarmRSI likewise describes experimental feedback guiding
mutation and recombination across scientific agents~\cite{swarmrsi}. Its
October 2026 preprint is a research report about an evolving feedback loop,
not evidence for a first-proposal, zero-feedback protocol.

\paragraph{Which objective, for which deployment?}
The intended ``Jev'' reference is the TypeSafe model introduced by Diogo
Almeida, not a researcher of that name~\cite{jev}. Its stated objective is to
return calibrated, typed decisions suitable for software integration. That
argument motivates matching a learning signal to the downstream task; it
does not establish that text generation is generally the wrong objective.
Human-preference learning and learning from verifiable outcomes are distinct,
already established approaches~\cite{rlhf,rlvr}, and contemporary language
models may combine several training signals. Exact training recipes for the
historical frontier entrants are not established by their submission names.
The appropriate comparison therefore concerns observed behavior, declared
information access, and measured resources, not a presumed dichotomy between
``preference-only'' language models and ``grounded'' algorithm models.

LeCun and coauthors argue that autonomous learning should integrate observation,
active behavior, and internal control over when to use each~\cite{lecun}.
This supplies a contextualized critique of current learning systems, not an
impossibility theorem about new knowledge. This manuscript avoids attributing
an unverified podcast quotation to LeCun. Its empirical question remains
answerable whichever broader architecture ultimately proves most capable:
after training has stopped, what useful algorithmic competence survives when
candidate feedback is withheld until commitment?
